\scanpage{157}
On l'appelle le \deleted{\unsure{secti|section}} \uline{jet de Picard-Fuchs défini par} $P'$ \deleted{\unsure{$\delta_P$}} \uline{et la stratification $\sigma$ donnée} \deleted{\unsure{de}} $P=P'\mathbin{\overset{E'}{\times}}E$. Ce n'est pas tt à fait un jet \inserted{intégrable} de $E''$, mais son \unsure{déguisement} \unsure{affine|affaibli} : c'est \unsure{là|le|à-dire} un tel jet, défini modulo \deleted{\unsure{une sec}} \inserted{l'image \inserted{par $vd^\infty_Eu$} d'une} section de \deleted{\unsure{$P^\infty_{S/T}(E')$}} $E'$.

Soit $F$ un module \inserted{qu. coh.} sur \deleted{$E$} $S$, et
\[
 D:E''\longrightarrow F \tag*{(39)}
\]
un op différentiel \deleted{\unsure{d'or}} rel. / $T$, définissant donc un hom \inserted{linéaire} continu : \deleted{\unsure{hom}}
\[
 u_D:P^\infty_{S/T}(E'')\longrightarrow F . \tag*{(40)}
\]
On dit que \deleted{$u_D$} $D$ est un \uline{opérateur de Picard-Fuchs}, \uline{relativement à la suite exacte} (33), si l'application linéaire correspondante (40) s'annule sur $v\delta(E)$, i.e. si $u_D$ se factorise en un homomorphisme
\[
 \overline u_D:P^\infty_{S/T}(E'')/v\delta(E)\longrightarrow F . \tag*{(41)}
\]
\deleted{Soit alors} Si alors $(P',\sigma)$ sont comme dessus, d'où un $\varphi$ (38), on peut former $\overline u_D(\varphi)\in\deleted{\unsure{H}}\Gamma(S,F)$ ; \unsure{on posera|en posant|on pose}
\[
 \deleted{\unsure{u_D(\varphi)=}}\ \mu_D(P,\sigma)=\overline u_D(\varphi(P,\sigma))\in\Gamma(S,F)\,, \tag*{(42)}
\]
% archive sheet number: 156
